\chapter{Discusión}
\label{ch:discusion}

Este capítulo interpreta los resultados experimentales a la luz de las hipótesis planteadas, analiza los mecanismos subyacentes que explican los patrones observados y transparenta las limitaciones del estudio con honestidad académica.

\section{Validación de las Hipótesis}

La evidencia experimental respalda sólidamente la hipótesis principal H1. TRSS supera a todas las referencias —tanto geométricas como GSP puramente espaciales— en todas las métricas cuando la pérdida de canales es contigua y severa. La magnitud del efecto es sustancial: ganancias de 5 a 10~dB en SNR y una reducción del error de latencia de componentes ERP en un factor de tres a cuatro veces. El mecanismo que explica esta superioridad no es una mejor representación espacial en sí misma, sino la capacidad de la información previa de naturaleza temporal para proveer información cuando la información espacial colapsa. En condiciones de pérdida contigua severa, los métodos espaciales se enfrentan a un problema fundamental: no hay suficientes electrodos buenos en la vecindad para interpolar con precisión, y el espacio de soluciones espacialmente suaves que son consistentes con las pocas observaciones disponibles es demasiado grande. TRSS reduce drásticamente este espacio al imponer una restricción adicional: la señal no solo debe ser espacialmente suave, sino también temporalmente continua. Esta restricción está anclada en la física de los generadores neuronales y resulta ser precisamente la información complementaria necesaria.

La hipótesis de equivalencia condicional H2 también recibe respaldo empírico. En condiciones de pérdida aleatoria igual o inferior al 20\% con montajes densos, MNE-Python es estadísticamente indistinguible de TRSS en varias métricas, y es superior en preservación espectral en algunos escenarios. Este resultado no debilita la contribución de la tesis; la fortalece, porque demuestra que el diseño experimental es lo suficientemente riguroso para revelar no solo cuándo un método es superior, sino también cuándo no lo es. Una evaluación comparativa honesta debe caracterizar los regímenes de equivalencia tanto como los de dominancia.

\section{Interpretación de los Mecanismos Subyacentes}

El estudio de ablación demostró que la combinación de los términos espacial y temporal produce una mejora mayor que la suma de sus efectos individuales, revelando una sinergia entre ambas fuentes de información previa. La interpretación fisiológica es que el grafo espacial le indica al algoritmo qué canales están relacionados —conectividad anatómica o funcional—, mientras que la penalización temporal le indica cómo debe evolucionar la señal en esos canales —con continuidad—. El grafo sin temporalidad produce superficies espacialmente suaves pero temporalmente ruidosas; la penalización temporal sin grafo produce señales suaves en el tiempo pero ignora qué canales son fisiológicamente adyacentes. Solo la combinación de ambas restricciones produce reconstrucciones que son simultáneamente coherentes en el espacio y en el tiempo.

El desempeño de MNE-Python podría parecer paradójico: ¿cómo puede un método formulado en 1989, sin información temporal ni topología de grafos, competir con algoritmos teóricamente superiores? La respuesta, documentada en detalle en el Capítulo~2, es que MNE-Python no es el algoritmo de 1989, sino el resultado de más de una década de mejoras incrementales de ingeniería que en conjunto producen un estimador sorprendentemente robusto. Ninguna de estas mejoras es revolucionaria por sí sola, pero su acumulación transforma un método teóricamente simple en una referencia formidable.

\section{Limitaciones del Estudio}

La transparencia sobre las limitaciones es esencial para la integridad académica. Este estudio asume que la pérdida de canales es permanente durante todo el segmento de análisis, sin evaluar escenarios de pérdida intermitente que son comunes en registros ambulatorios prolongados. Se asumió una referencia promedio común para todos los métodos, sin variar sistemáticamente la elección de la referencia, la cual afecta la distribución espacial del potencial. Los canales se eliminaron artificialmente de datos completos, mientras que en la práctica clínica los canales malos frecuentemente contienen artefactos —no solo ausencia de señal—, lo que podría afectar diferencialmente a los métodos que dependen de los canales vecinos. La optimización de hiperparámetros se realizó sobre un subconjunto de los mismos conjuntos de datos usados para la evaluación final, cuando idealmente debería haberse realizado sobre conjuntos completamente independientes. No se evaluó el impacto en tareas subsecuentes como la clasificación en interfaces cerebro-computador o la localización de fuentes. Finalmente, TRSS es entre 50 y 100 veces más lento que MNE-Python, una diferencia que puede ser prohibitiva para aplicaciones con restricciones de tiempo real.

\section{Implicaciones para la Práctica}

Los resultados de esta tesis permiten formular una guía práctica para la selección del método de interpolación. Se recomienda utilizar MNE-Python cuando el montaje es denso —sesenta canales o más—, la pérdida de canales es igual o inferior al 20\% y mayoritariamente aleatoria, se requiere procesamiento en tiempo real o casi real, la preservación espectral es prioritaria y la morfología fina de los ERPs no es crítica, o se busca simplicidad de implementación con mínima dependencia de bibliotecas externas. Por el contrario, se recomienda utilizar TRSS cuando la pérdida de canales es igual o superior al 30\% y contigua —bloques de amplificadores fallando—, la preservación de la morfología de ERPs es crítica —estudios de potenciales evocados—, se requiere la máxima fidelidad de reconstrucción sin restricciones estrictas de tiempo, o el montaje es disperso —32 canales o menos—, donde cada canal tiene un alto valor informativo.

Para la comunidad de investigación en GSP, este trabajo ofrece tres lecciones. La primera es que los métodos de referencia deben optimizarse tanto como el método propuesto: la diferencia entre una implementación ingenua y MNE-Python es de 3 a 8~dB en SNR. La segunda es que deben reportarse los regímenes de falla, no solo los de éxito: un método que gana en ocho de diez escenarios pero pierde catastróficamente en los otros dos es menos útil que uno que empata en los diez. La tercera es que la regularización temporal es el factor diferencial, no la topología del grafo en sí misma, lo que sugiere que las mejoras futuras deberían enfocarse en refinar el modelo temporal más que en optimizar la construcción del grafo.